\centerline{\bf \Huge Делимость и остатки}

\begin{flushright}
        \scriptsize{Yare yare daze}
\end{flushright}

В музее старых денег стоит чудесный разменный автомат «старина MOD3», действующий следующим образом. Если в него опустить несколько монет, то в обмен он вам выдаст ту же сумму денег, но копеечными монетами старого образца: вначале он будет отсчитывать трехкопеечные монеты до тех пор, пока сумма, которую ему осталось выплатить, не станет меньше 3. После этого он выдаст сдачу - остаток. Ясно, что этим остатком может быть любое из трех чисел: $0,1,2$ (т.е. если исходное число не делилось на три, то автомат выдаст вам помимо трехкопеечных монет еще одну копеечную или двухкопеечную монетку).


\problem{1}
Что выдаст автомат, если в него опустить монету 1 рубль? А если одновременно опустить две монеты: рублёвую и пятирублёвую?

\problem{2}
Вам повезло: вы попали в этот музей и у вас есть три пятирублёвые монеты. Удастся ли вам заполучить старинные монеты всех трех видов?

\problem{3}
Вы опустили в чудо-автомат одновременно 22 пятирублёвые монеты и 33 рублёвые. Какую сдачу вы получите?

\problem{4}
Рядом есть автомат «гуманоид MOD7». Oн действует так же, как и «старина MOD3», но отсчитывает вначале 7 -копеечные монеты, а потом выдает сдачу остаток. Какую сдачу вы получите, опустив в него одновременно 22 рублёвые монеты и 33 пятирублёвые?

\problem{5}
Какую сдачу можно получить, опустив в автомат «гуманоид MOD7» одновременно $7 k+20$ монет по 5 рублей?

\problem{6}
Какую сдачу можно получить, если опустить в «гуманоид MOD7» одновременно $7 k+20$ монет по $7 t+5$ рублей каждая?

\problem{7}
Если опустить некоторую монету в автомат «MOD12», то он выдаст сдачу 6 коп. А какую сдачу выдаст автомат «MOD4», если опустить в него эту же монету?

\problem{8}
Если опустить некоторую монету в автомат «старина MOD3», то на сдачу он выдаст 2 коп. Если эту же монету опустить в автомат «MOD4», то сдачи не будет. Какую сдачу выдал бы автомат «MOD12», если бы эту монету опустили в него?